% Generated by tools/generate_statements_only_paper.py.
% Do not edit this derived file; edit the canonical paper source instead.

\section{Structural Reduction to Finite Cases}
\label{sec:structural-reductions}

Throughout, indices are taken in $\mathbb Z/6\mathbb Z$.  Put
\[
 O=(0,0),\qquad
 V_i=\left(\cos\frac{i\pi}{3},\sin\frac{i\pi}{3}\right),
 \qquad
 H=\mathrm{conv}\{V_0,\ldots,V_5\}.
\]
Thus $H$ is the regular hexagon of side length one.  We use the notation
\[
 e_{i,i+1}=[V_i,V_{i+1}],\qquad
 r_i=[O,V_i],\qquad M_i=\frac12V_i,
\]
and, when needed,
\[
 D=\bigcup_{i=0}^5r_i,\qquad
 S=\partial H\cup D,\qquad
 S_{1/2}=\partial H\cup\{O,M_0,\ldots,M_5\}.
\]
The arms $r_i$ meet only at $O$, up to sets of one-dimensional measure zero,
and consequently
\begin{equation}
 \mathcal H^1(\partial H)=6,
 \qquad \mathcal H^1(D)=6,
 \qquad \mathcal H^1(S)=12.
 \label{eq:target-lengths}
\end{equation}
For $L>0$, $H_L:=LH$ denotes the concentric regular hexagon of side length
$L$.  All symmetries used below belong to the dihedral group $D_6$; applying
a symmetry always transports the indices, orientations, and incoming and
outgoing labels together.

\subsection{Open, closed, and scaled formulations}

\begin{proposition}[Compactness and scaling equivalence]
\label{prop:open-closed-scaled}
The following assertions are equivalent.
\begin{enumerate}
 \item The set $H$ is covered by seven open equilateral triangles of side
 length $1$.
 \item For some $\varepsilon>0$, the set $H$ is covered by seven closed
 equilateral triangles of side length $1-\varepsilon$.
 \item For some $L>1$, the set $H_L$ is covered by seven closed equilateral
 triangles of side length $1$.
\end{enumerate}
\end{proposition}

% Proof environment omitted (source: 02_structural_reductions.tex).

We conduct the geometric argument in the open side-one model.  Denote the
original open triangles by $U_C$ and $U_i$ for $0\le i\le5$, and put
\[
 T_C=\overline{U_C},\qquad T_i=\overline{U_i}.
\]
The distinction between $U_i$ and $T_i$ will be maintained: closure may add
endpoints or tangencies, although it does not change area or one-dimensional
trace measure.

\begin{lemma}[Distinct triangles]
\label{lem:distinct-roles}
The seven triangles can be labelled so that
\[
 O\in U_C,\qquad V_i\in U_i\quad(0\le i\le5),
\]
and these seven triangles are distinct.  Consequently
\[
 O\in\mathrm{int}(T_C),\qquad V_i\in\mathrm{int}(T_i).
\]
\end{lemma}

% Proof environment omitted (source: 02_structural_reductions.tex).

\subsection{The center classification}

For a closed C triangle define
\[
 N_E(T_C)=
 \left\lvert\left\{i:T_C\cap e_{i,i+1}\text{ contains a
 positive-length interval}\right\}\right\rvert.
\]
Point contacts, endpoint contacts, and tangencies do not contribute.  We say
that $T_C$ has type CE0, CE1, or CE2 according as $N_E(T_C)$ is $0$, $1$, or
$2$.

\begin{proposition}[Exhaustive center types]
\label{prop:ce-classification}
Every C triangle has exactly one of the types CE0, CE1, and CE2.
Equivalently,
\[
 N_E(T_C)\le2.
\]
\end{proposition}

% Proof environment omitted (source: 02_structural_reductions.tex).

For the CE1/CE2 types, Proposition~\ref{prop:unique-center-midpoint} below
records the additional fact that $T_C$ contains exactly one radial midpoint.

\subsection{Vertex types and the T3-like normalization}

For the closure $T_i$ of an original open V triangle $U_i$, let
\[
 o(T_i)=\left\lvert\left\{
 \text{vertices of }T_i\text{ outside }H
 \right\}\right\rvert
\]
and put
\[
 n(T_i)=\left\lvert\left\{j\in\{i-1,i+1\}:
 \Haus(T_i\cap r_j)>0\right\}\right\rvert.
\]
We use the names
\[
 \begin{array}{c|c}
 \text{type}&(o,n)\\ \hline
 \mathrm{Vd0}&(1,0),(2,0),(3,0)\\
 \mathrm{Vd1}&(1,1)\\
 \mathrm{Vd2}&(1,2)\\
 \text{T3-like}&(2,1).
 \end{array}
\]

\begin{lemma}[Local wedge]
\label{lem:local-wedge}
Every original V triangle has $1\le o\le3$.  If
$\Haus(T_i\cap r_j)>0$ for some $j\in\{i-1,i+1\}$, at least one of its
vertices belongs to $H$; if both $\Haus(T_i\cap r_{i-1})>0$ and
$\Haus(T_i\cap r_{i+1})>0$, at least two of its vertices belong to $H$.
\end{lemma}

% Proof environment omitted (source: 02_structural_reductions.tex).

\begin{proposition}[Exhaustive vertex types]
\label{prop:vertex-classification}
Every original V triangle has exactly one of the four types Vd0, Vd1, Vd2,
and T3-like.
\end{proposition}

% Proof environment omitted (source: 02_structural_reductions.tex).

The next normalization is useful for trace estimates, but its final warning
is essential.

\begin{proposition}[T3-like closed-trace domination]
\label{prop:t3-translation}
Let $T$ be the closure of an original T3-like $V_i$-triangle.  There is a
translate $\widehat T$ of $T$, with the same orientation and side length,
such that
\[
 V_i\in\mathrm{relint}(\text{a side of }\widehat T),
 \qquad T\cap H\subsetneq\widehat T\cap H,
\]
and $\widehat T$ remains of type $(o,n)=(2,1)$.  Every old open interior
point in $H\setminus\{V_i\}$ remains an interior point of $\widehat T$.
The point $V_i$ itself becomes a boundary point, so $\widehat T$ is a
closed-trace majorant and is not a replacement open triangle.
\end{proposition}

% Proof environment omitted (source: 02_structural_reductions.tex).

\subsection{Reach coordinates, supercritical V triangles, and gaps}

For the original V triangle $U_i$, define its maximal incoming,
outgoing, and radial reaches by
\[
 \begin{aligned}
 A(U_i)&=
 \sup_{\substack{t\in[0,1]\\V_i+t(V_{i-1}-V_i)\in U_i}}t,\\
 B(U_i)&=
 \sup_{\substack{t\in[0,1]\\V_i+t(V_{i+1}-V_i)\in U_i}}t,\\
 C(U_i)&=
 \sup_{\substack{t\in[0,1]\\V_i+t(O-V_i)\in U_i}}t.
 \end{aligned}
\]
Closure does not change the suprema, so
\[
 A(U_i)=A(T_i),\qquad B(U_i)=B(T_i),\qquad C(U_i)=C(T_i).
\]
For compact indexed formulas, put
\[
 A_i:=A(U_i)=A(T_i),\qquad
 B_i:=B(U_i)=B(T_i),\qquad
 C_i:=C(U_i)=C(T_i).
\]
The closure $T_i$ contains the corresponding endpoints. We reserve
lowercase $(a_i,b_i,c_i)$ for selected or prescribed reach lower bounds
whenever actual and selected values occur together.  A realizing triangle then
satisfies
\[
 A_i\ge a_i,\qquad B_i\ge b_i,\qquad C_i\ge c_i.
\]
An actual V triangle is supercritical if $A_i+B_i>1$, and
\begin{equation}
 N_+=\left\lvert\left\{i:A_i+B_i>1\right\}\right\rvert.
 \label{eq:N-plus-definition}
\end{equation}
This definition never counts a smaller reach pair in place of the actual
V triangle.

Parametrize $e_{i,i+1}$ by
\[
 X_i(x)=V_i+x(V_{i+1}-V_i),\qquad 0\le x\le1.
\]
The two incident open traces are
\[
 U_i\cap e_{i,i+1}=[0,B_i),
 \qquad
 U_{i+1}\cap e_{i,i+1}=(1-A_{i+1},1]
\]
in this coordinate.  A \emph{boundary gap} on this edge is the nonempty
closed complement between these two traces,
\[
 [B_i,1-A_{i+1}]\qquad\text{when }B_i\le1-A_{i+1}.
\]
In particular, equality gives a singleton gap.  This convention retains the
actual open traces; a point missing from both open triangles is not erased merely
because both closures contain it.

\begin{lemma}[Exhaustiveness of the gap split]
\label{lem:gap-exhaustion}
On each edge the complement of the two incident V triangle traces is either
empty or one closed interval; a singleton interval counts as a gap.  Every
edge having a gap has positive-length intersection with the open center
triangle.  Consequently a $\CEzero$ center permits no gaps, a $\CEone$ center
permits at most one, and a $\CEtwo$ center permits at most two.
\end{lemma}

% Proof environment omitted (source: 02_structural_reductions.tex).

\begin{lemma}[T3-like V triangles are nonsupercritical]
\label{lem:t3-nonsupercritical}
Every original T3-like V triangle satisfies
\begin{equation}
 A_i+B_i\le1.
\label{eq:t3-nonsupercritical}
\end{equation}
In particular, it is not counted by $N_+$.
\end{lemma}

% Proof environment omitted (source: 02_structural_reductions.tex).

We state the precise passage from maximal reaches to strict handoff lower bounds.
Its hypothesis that the six V triangles themselves cover the boundary is
part of the statement.

\begin{proposition}[Strict boundary handoffs]
\label{prop:strict-handoffs}
Assume that $U_0,\ldots,U_5$ cover $\partial H$.  On each edge choose
\[
 \xi_i\in(1-A_{i+1},B_i)
\]
and define
\[
 a_i=1-\xi_{i-1},\qquad b_i=\xi_i.
\]
Then $0<a_i<A_i$, $0<b_i<B_i$, the two selected anchors lie in $U_i$, and
\[
 a_i^2+a_ib_i+b_i^2<1.
\]
Moreover:
\begin{enumerate}
 \item if exactly one actual V triangle, indexed by $p$, is supercritical, every
 strict selection has exactly one selected supercritical V triangle, also indexed
 by $p$;
 \item if at least two actual V triangles are supercritical, some simultaneous
 strict selection has at least two selected supercritical V triangles.
\end{enumerate}
\end{proposition}

% Proof environment omitted (source: 02_structural_reductions.tex).

\subsection{Midpoint forcing}

We first need a local fact for V triangles.  Take unit vectors
\[
 u=\left(\frac12,\frac{\sqrt3}{2}\right),\qquad
 v=\left(\frac12,-\frac{\sqrt3}{2}\right).
\]

\begin{lemma}[Self-midpoint obstruction]
\label{lem:self-midpoint}
Let a closed unit triangle contain
\[
 0,\qquad au,\qquad bv,\qquad \frac{u+v}{2}.
\]
Then $a+b\le1$.  Consequently, for every V triangle,
\[
 M_i\in T_i\quad\Longrightarrow\quad A_i+B_i\le1.
\]
If the two demanded boundary points and the midpoint are contained with a
positive margin, the conclusion is strict.
\end{lemma}

% Proof environment omitted (source: 02_structural_reductions.tex).

The corresponding center fact is exact and uses the interior condition at
$O$.

\begin{proposition}[CE1/CE2 exactly-one-midpoint property]
\label{prop:unique-center-midpoint}
Let $T_C$ be a closed unit equilateral triangle with
$O\in\mathrm{int}(T_C)$.  If $T_C$ is CE1 or CE2, equivalently if it has
positive-length intersection with a boundary edge of $H$, then
\[
 \left\lvert T_C\cap\{M_0,\ldots,M_5\}\right\rvert=1.
\]
\end{proposition}

% Proof environment omitted (source: 02_structural_reductions.tex).

\begin{remark}
The hypothesis $O\in\mathrm{int}(T_C)$ cannot be weakened to
$O\in T_C$: the triangle $\mathrm{conv}\{O,V_0,V_1\}$ has a boundary-edge
overlap and contains both $M_0$ and $M_1$.
\end{remark}

\begin{proposition}[Exhaustive structural reduction]
\label{prop:exhaustive-structural-reduction}
Every hypothetical cover by seven open unit equilateral triangles admits the
distinct center and V triangles used in Table~\ref{tab:routing}, and its
maximal reaches determine exactly one row of that table.
\end{proposition}

% Proof environment omitted (source: 02_structural_reductions.tex).
